\documentclass[11pt,a4paper]{article}
\usepackage[utf8]{inputenc}
\usepackage{amsmath,amssymb,amsfonts}
\usepackage[margin=2cm]{geometry}
\usepackage{booktabs}
\usepackage{longtable}
\usepackage{array}
\usepackage{hyperref}
\usepackage{xcolor}

\definecolor{pass}{HTML}{2ea043}
\definecolor{fail}{HTML}{cf222e}
\definecolor{warn}{HTML}{d4880a}

\newcommand{\Iref}{I^{\mathrm{ref}}}
\newcommand{\Ifill}{I^{\mathrm{ref,filled}}}
\newcommand{\IDelta}{I_{\Delta}}
\let\oldcheckmark\checkmark
\renewcommand{\checkmark}{\textcolor{pass}{\bfseries\oldcheckmark}}
\newcommand{\crossmark}{\textcolor{fail}{\bfseries\sffamily X}}
\newcommand{\passmark}{\textcolor{pass}{PASS}}
\newcommand{\failmark}{\textcolor{fail}{FAIL}}

\begin{document}

\title{Refined 3D Index — Full Report \\ \large 5\_1}
\author{Generated by \texttt{manifold-index}}
\date{\today}
\maketitle
\tableofcontents
\newpage

\section{Manifold Overview}
Manifold name: \textbf{5\_1}

\begin{itemize}
  \item Number of ideal tetrahedra: $n = 3$
  \item Number of cusps: $r = 1$
  \item Internal edges: $n - r = 2$ (easy: 1, hard: 1)
  \item Computation order: $q^{24/2}$
  \item Boundary curve range: $m \in \{-2, -1, 0, 1, 2\}$, $e \in \{-2, -\tfrac{3}{2}, -1, -\tfrac{1}{2}, 0, \tfrac{1}{2}, 1, \tfrac{3}{2}, 2\}$
\end{itemize}

\section{Ideal Triangulation (SnaPy Data)}
Gluing equation matrix from SnaPy, shape $(5 \times 9)$.

\begin{longtable}{c c c c}
\toprule
Edge & $(Z_{1}, Z_{1}', Z_{1}'')$ & $(Z_{2}, Z_{2}', Z_{2}'')$ & $(Z_{3}, Z_{3}', Z_{3}'')$ \\ \midrule
\endhead
0 & $(1, 0, 1)$ & $(1, 2, 0)$ & $(1, 1, 0)$ \\
1 & $(0, 2, 1)$ & $(1, 0, 2)$ & $(0, 1, 2)$ \\
2 & $(1, 0, 0)$ & $(0, 0, 0)$ & $(1, 0, 0)$ \\
\bottomrule
\end{longtable}

\begin{longtable}{c c c c}
\toprule
Cusp & $(Z_{1}, Z_{1}', Z_{1}'')$ & $(Z_{2}, Z_{2}', Z_{2}'')$ & $(Z_{3}, Z_{3}', Z_{3}'')$ \\ \midrule
\endhead
$\alpha$ & $(0, 0, -1)$ & $(0, -1, 0)$ & $(0, 0, 1)$ \\
$\beta$ & $(1, 1, -7)$ & $(0, -10, 0)$ & $(-1, 1, 9)$ \\
\bottomrule
\end{longtable}

\section{Edge Classification}
Easy edges: 1 total, 1 independent. Hard edges: 1.

\begin{longtable}{c l l}
\toprule
\multicolumn{3}{l}{\textbf{Hard Edges}} \\ \midrule
Index & Symbolic Equation & Triplets $(f_i, g_i, h_i)$ \\ \midrule
\endhead
H$_{0}$ & $2Z_{1}' + Z_{1}'' + Z_{2} + 2Z_{2}'' + Z_{3}' + 2Z_{3}'' = 2$ & $(0,2,1)\;(1,0,2)\;(0,1,2)$ \\
\bottomrule
\end{longtable}

\begin{longtable}{c c l l}
\toprule
\multicolumn{4}{l}{\textbf{Easy Edges}} \\ \midrule
Index & Independent? & Symbolic Equation & Triplets $(f_i, g_i, h_i)$ \\ \midrule
\endhead
E$_{0}$ & \checkmark & $Z_{1} + Z_{3} = 2$ & $(1,0,0)\;(0,0,0)\;(1,0,0)$ \\
\bottomrule
\end{longtable}

\section{Neumann--Zagier Data}
$g_{\text{NZ}} \in \mathrm{Sp}(2n, \mathbb{Q})$, $6 \times 6$, column ordering $(Z_1, \ldots, Z_n, Z_1'', \ldots, Z_n'')$.
{\small
\[
g_{\text{NZ}} = \left(\begin{array}{rrrrrr}
  0 & 1 & 0 & -1 & 1 & 1 \\  %\; \mu_{0}
  -2 & 1 & -1 & -1 & 2 & 1 \\  %\; \text{H}_{0}
  1 & 0 & 1 & 0 & 0 & 0 \\  %\; \text{E}_{0}
  0 & 5 & -1 & -4 & 5 & 4 \\  %\; \lambda_{0}/2
  0 & 0 & 1 & -1 & 1 & 1 \\  %\; \Gamma_{0}
  -2 & 0 & 0 & -1 & 2 & 2 \\  %\; \Gamma_{1}
\end{array}\right)
\]
}

Row labels (right margin): 
$\mu_{0}$, $\text{H}_{0}$, $\text{E}_{0}$, $\lambda_{0}/2$, $\Gamma_{0}$, $\Gamma_{1}$.

Affine shifts: $\nu_x = (-1, 1, -2)$, $\nu_p = (-4, 0, 0)$.

\section{3D Index Formula}
The 3D index is computed via:
\[
  I(\mathbf{m}_{\text{ext}}, \mathbf{e}_{\text{ext}}) = \sum_{\mathbf{e}_{\text{int}} \in (\tfrac{1}{2}\mathbb{Z})^{n-r}} (-q^{1/2})^{\mathbf{m} \cdot \nu_p \;-\; \mathbf{e} \cdot \nu_x} \;\prod_{a=1}^{n} \IDelta(m_a, e_a)
\]
where $\kappa = (\mathbf{m}_{\text{ext}},\, \mathbf{0}^{n-r},\, \mathbf{e}_{\text{ext}},\, \mathbf{e}_{\text{int}})$ and the local charges are:
\[
  m_a = (g_{\text{NZ}}^{-1}\kappa)_a, \qquad   e_a = (g_{\text{NZ}}^{-1}\kappa)_{n+a}
\]
Only terms with integer $(m_a, e_a)$ contribute.

\noindent\textbf{Explicit Local Charges.}
$\Delta_{1}$: $\;m_{1} = -4m_{0} + e_{0} + e_{\text{int},0}$, $\;e_{1} = -2e_{\text{int},0} + e_{\text{int},1}$

$\Delta_{2}$: $\;m_{2} = 5m_{0} - e_{0} - 2e_{\text{int},0}$, $\;e_{2} = -5m_{0} + e_{0} + e_{\text{int},0}$

$\Delta_{3}$: $\;m_{3} = 4m_{0} - e_{0} - e_{\text{int},0}$, $\;e_{3} = m_{0} - e_{\text{int},0} + e_{\text{int},1}$

The tetrahedron index $\IDelta(m, e)$ is given by
\[
  I_t(m, e) = \sum_{k=\max(0,-e)}^{\infty} \frac{(-1)^k \, q^{k(k+1)/2 - (2k+e)m/2}}{\prod_{j=1}^{k}(1-q^j) \,\prod_{j=1}^{k+e}(1-q^j)}
\]
with $\IDelta(m,e) = (-q^{1/2})^m \cdot I_t(-m-e, m)$ when $m + e \ge 0$, else $\IDelta(m,e) = I_t(m, e)$.

\section{Refined Index Results}
Series truncated at $q^{24/2}$ (i.e.\ $N_{\max} = 12$).
Number of computed sectors: 45.
Non-zero sectors: 1.

\begin{align*}
  \Iref(\mathbf{m}=(-2),\;\mathbf{e}=(-2)) &= 0 \\
  \Iref(\mathbf{m}=(-2),\;\mathbf{e}=(-3/2)) &= 0 \\
  \Iref(\mathbf{m}=(-2),\;\mathbf{e}=(-1)) &= 0 \\
  \Iref(\mathbf{m}=(-2),\;\mathbf{e}=(-1/2)) &= 0 \\
  \Iref(\mathbf{m}=(-2),\;\mathbf{e}=(0)) &= 0 \\
  \Iref(\mathbf{m}=(-2),\;\mathbf{e}=(1/2)) &= 0 \\
  \Iref(\mathbf{m}=(-2),\;\mathbf{e}=(1)) &= 0 \\
  \Iref(\mathbf{m}=(-2),\;\mathbf{e}=(3/2)) &= 0 \\
  \Iref(\mathbf{m}=(-2),\;\mathbf{e}=(2)) &= 0 \\
  \Iref(\mathbf{m}=(-1),\;\mathbf{e}=(-2)) &= 0 \\
  \Iref(\mathbf{m}=(-1),\;\mathbf{e}=(-3/2)) &= 0 \\
  \Iref(\mathbf{m}=(-1),\;\mathbf{e}=(-1)) &= 0 \\
  \Iref(\mathbf{m}=(-1),\;\mathbf{e}=(-1/2)) &= 0 \\
  \Iref(\mathbf{m}=(-1),\;\mathbf{e}=(0)) &= 0 \\
  \Iref(\mathbf{m}=(-1),\;\mathbf{e}=(1/2)) &= 0 \\
  \Iref(\mathbf{m}=(-1),\;\mathbf{e}=(1)) &= 0 \\
  \Iref(\mathbf{m}=(-1),\;\mathbf{e}=(3/2)) &= 0 \\
  \Iref(\mathbf{m}=(-1),\;\mathbf{e}=(2)) &= 0 \\
  \Iref(\mathbf{m}=(0),\;\mathbf{e}=(-2)) &= 0 \\
  \Iref(\mathbf{m}=(0),\;\mathbf{e}=(-3/2)) &= 0 \\
  \Iref(\mathbf{m}=(0),\;\mathbf{e}=(-1)) &= 0 \\
  \Iref(\mathbf{m}=(0),\;\mathbf{e}=(-1/2)) &= 0 \\
  \Iref(\mathbf{m}=(0),\;\mathbf{e}=(0)) &= 1 + q \, \eta_0^{-1} -q + q^{2} \, \eta_0^{-1} -2 \, q^{2} + q^{2} \, \eta_0 + q^{3} \, \eta_0^{-1} -2 \, q^{3} + q^{3} \, \eta_0 + q^{4} \, \eta_0^{-2} -2 \, q^{4} + q^{4} \, \eta_0 + q^{5} \, \eta_0^{-2} -q^{5} \, \eta_0^{-1} + 2 \, q^{6} \, \eta_0^{-2} -3 \, q^{6} \, \eta_0^{-1} + q^{6} -q^{6} \, \eta_0 + q^{6} \, \eta_0^{2} + 2 \, q^{7} \, \eta_0^{-2} -5 \, q^{7} \, \eta_0^{-1} + 5 \, q^{7} -3 \, q^{7} \, \eta_0 + q^{7} \, \eta_0^{2} + 3 \, q^{8} \, \eta_0^{-2} -7 \, q^{8} \, \eta_0^{-1} + 7 \, q^{8} -5 \, q^{8} \, \eta_0 + 2 \, q^{8} \, \eta_0^{2} + q^{9} \, \eta_0^{-3} + 2 \, q^{9} \, \eta_0^{-2} -9 \, q^{9} \, \eta_0^{-1} + 11 \, q^{9} -7 \, q^{9} \, \eta_0 + 2 \, q^{9} \, \eta_0^{2} + q^{10} \, \eta_0^{-3} + 2 \, q^{10} \, \eta_0^{-2} -10 \, q^{10} \, \eta_0^{-1} + 13 \, q^{10} -9 \, q^{10} \, \eta_0 + 3 \, q^{10} \, \eta_0^{2} + 2 \, q^{11} \, \eta_0^{-3} -10 \, q^{11} \, \eta_0^{-1} + 16 \, q^{11} -10 \, q^{11} \, \eta_0 + 2 \, q^{11} \, \eta_0^{2} + 3 \, q^{12} \, \eta_0^{-3} -2 \, q^{12} \, \eta_0^{-2} -8 \, q^{12} \, \eta_0^{-1} + 14 \, q^{12} -10 \, q^{12} \, \eta_0 + 2 \, q^{12} \, \eta_0^{2} + q^{12} \, \eta_0^{3} \\
  \Iref(\mathbf{m}=(0),\;\mathbf{e}=(1/2)) &= 0 \\
  \Iref(\mathbf{m}=(0),\;\mathbf{e}=(1)) &= 0 \\
  \Iref(\mathbf{m}=(0),\;\mathbf{e}=(3/2)) &= 0 \\
  \Iref(\mathbf{m}=(0),\;\mathbf{e}=(2)) &= 0 \\
  \Iref(\mathbf{m}=(1),\;\mathbf{e}=(-2)) &= 0 \\
  \Iref(\mathbf{m}=(1),\;\mathbf{e}=(-3/2)) &= 0 \\
  \Iref(\mathbf{m}=(1),\;\mathbf{e}=(-1)) &= 0 \\
  \Iref(\mathbf{m}=(1),\;\mathbf{e}=(-1/2)) &= 0 \\
  \Iref(\mathbf{m}=(1),\;\mathbf{e}=(0)) &= 0 \\
  \Iref(\mathbf{m}=(1),\;\mathbf{e}=(1/2)) &= 0 \\
  \Iref(\mathbf{m}=(1),\;\mathbf{e}=(1)) &= 0 \\
  \Iref(\mathbf{m}=(1),\;\mathbf{e}=(3/2)) &= 0 \\
  \Iref(\mathbf{m}=(1),\;\mathbf{e}=(2)) &= 0 \\
  \Iref(\mathbf{m}=(2),\;\mathbf{e}=(-2)) &= 0 \\
  \Iref(\mathbf{m}=(2),\;\mathbf{e}=(-3/2)) &= 0 \\
  \Iref(\mathbf{m}=(2),\;\mathbf{e}=(-1)) &= 0 \\
  \Iref(\mathbf{m}=(2),\;\mathbf{e}=(-1/2)) &= 0 \\
  \Iref(\mathbf{m}=(2),\;\mathbf{e}=(0)) &= 0 \\
  \Iref(\mathbf{m}=(2),\;\mathbf{e}=(1/2)) &= 0 \\
  \Iref(\mathbf{m}=(2),\;\mathbf{e}=(1)) &= 0 \\
  \Iref(\mathbf{m}=(2),\;\mathbf{e}=(3/2)) &= 0 \\
  \Iref(\mathbf{m}=(2),\;\mathbf{e}=(2)) &= 0
\end{align*}

\section{Weyl Symmetry}
Define the Weyl-shifted index:
\[
  f(\eta_j; m, e) = \eta_j^{a_j \cdot e + b_j \cdot m} \cdot \Iref(m, e)
\]
Weyl symmetry requires $f(m, e) = f(-m, -e)$ for each sector pair, where $f$ denotes the Weyl-shifted index above.

Weyl symmetry $f(m,e)=f(-m,-e)$: 0/45 sectors pass \failmark
 (45 failures may be due to missing partner sectors or finite truncation at large $|m|$, $|e|$).

The adjoint $\mathrm{su}(2)$ projection (eq.~2.59--2.61) integrates out $\eta$ and cusp fugacities against the Haar measure weighted by the adjoint character; the result must equal $-1$ per filled cusp.

Intermediate $c_e$ values ($(q^1, \eta^0)$ coefficient of $\Iref(m=0, e)$):
\begin{center}
\begin{tabular}{c c}
\toprule
$e$ & $c_e$ \\ \midrule
  $-2$ & $0$ \\
  $-\tfrac{3}{2}$ & $0$ \\
  $-1$ & $0$ \\
  $-\tfrac{1}{2}$ & $0$ \\
  $0$ & $-1$ \\
  $\tfrac{1}{2}$ & $0$ \\
  $1$ & $0$ \\
  $\tfrac{3}{2}$ & $0$ \\
  $2$ & $0$ \\
\bottomrule
\end{tabular}
\end{center}

Projected value: $0$ (expected $-1$): \failmark


\end{document}
